% !PRINTED 249
\weberpage{249}
%<<CONT
rational von $j(\omega)$ abhängen. Diese Gleichungen nennen wir
\gesperrt{Multiplikatorgleichungen erster Stufe}\footnote{Diese Gleichungen sind besonders eingehend von \gesperrt{Kiepert} untersucht worden, zuerst in mehreren Abhandlungen in Crelles Journal, Bd.~87, 88, 95, am ausführlichsten in den Abhandlungen in den Mathematischen Annalen, Bd.~26, 33. Vgl. auch F. \gesperrt{Klein}, Mathematische Annalen, Bd.~14, 15 und die oben erwähnten autographierten Vorlesungen.}.

% [sic] The source cites § 68, (16); the displayed quantity is instead the final unnumbered formula derived after (18).
Wir betrachten die Größen [§\,68, (16)]:

\weq{1}{P_{c,\eth,a} = i^{\frac{a-1}{2}}\left(\frac{c}{e}\right)\sqrt{\eth}\,\frac{\eta\left(\dfrac{c+\eth\,\omega}{a}\right)}{\eta(\omega)},}

und den Einfluss, den die Transformationen
\[
\begin{pmatrix} 1, & 0 \\ 1, & 1 \end{pmatrix}, \quad
\begin{pmatrix} 0, & 1 \\ -1, & 0 \end{pmatrix}
\]
auf sie ausüben. Dieser Einfluss bestimmt sich nach den Formeln
(8) bis (13), §\,69, wonach [weil $c_1 \equiv c \pmod{e}$]
\[
P_{c,\eth,a} \text{ durch } (\omega, \omega+1) \text{ in } e^{\frac{\pi i(\lambda-1)}{12}} P_{c_1,\eth,a}
\]

% Source scan confirms that the matrix and following fraction share one E(...) delimiter pair.
\weq{2}{\text{durch } \left(\omega, -\frac{1}{\omega}\right) \text{ in } \sqrt{\frac{\eth}{\eth_2}}\,\frac{i^{\frac{a-a_2}{2}}}{\sqrt{-i\omega}}\left(\frac{c}{e}\right)\left(\frac{c_2}{e_2}\right) E\!\left(\begin{matrix}\alpha, & \beta \\ \gamma, & \delta' \end{matrix}\ \frac{c_2+\eth_2\,\omega}{a_2}\right) P_{c_2,\eth_2,a_2}}

übergeht, worin $E$ die in §\,38, (15) angegebene Bedeutung hat.
Es ist aber [§\,69, (11), (12)]
\[
\sqrt{\eth}\,\sqrt{\alpha+\beta\,\frac{c_2+\eth_2\,\omega}{a_2}} = \sqrt{\eth_2\,\omega},
\]
und mithin, wenn wir
\[
E\!\left(\begin{matrix}\alpha, & \beta \\ \gamma, & \delta' \end{matrix}\ \frac{c_2+\eth_2\,\omega}{a_2}\right) = r\,\sqrt{-i\omega}\,\sqrt{\frac{\eth_2}{\eth}}
\]
setzen,
\[
r = \left(\frac{\alpha}{\beta}\right) i^{\frac{1-\beta}{2}} e^{\frac{\pi i}{12}\left[\beta(\alpha+\delta)-(\beta^2-1)\alpha\gamma\right]}
\]
\gesperrt{eine 24ste Einheitswurzel}, und durch die Vertauschung
\[
\left(\omega, -\frac{1}{\omega}\right)
\]
geht
\[
P_{c,\eth,a} \text{ in } r\,i^{\frac{a-a_2}{2}}\left(\frac{c}{e}\right)\left(\frac{c_2}{e_2}\right) P_{c_2,\eth_2,a_2}
\]
über.
